\begin{solapartat}
Per fer cada punt la màquina de cosir ha de fer una oscil·lació completa, per tant tindrem:
\[ \nu = \frac{1800\text{punts}}{60\mathrm{s}} = 30\ \mathrm{Hz} \ \text{\punts{0,3}} \]
Dues vegades l'amplitud del moviment serà igual a 20~mm $\Rightarrow$ A $=$ 10~mm $=$ 0,01~m \punts{0,3} L'equació del moviment serà:
\[ y(t) = A\,\cos(2\pi\nu t) \ \text{\punts{0,2}} = 0,01\,\cos(1,88\cdot 10^{2}t)\ \mathrm{m} \ \text{\punts{0,2}} \]
En el cas de que fessin servir la funció sinus enlloc del cosinus haurien de posar-hi una fase addicional de $\frac{\pi}{2}$
\end{solapartat}

\begin{solapartat}
\begin{align*}
v_y(t) &= \frac{\mathrm{d}y}{\mathrm{d}t} = -2\pi\nu A\sin(2\pi\nu t) \ \text{\punts{0,3}}\\
&\Rightarrow v_y(\text{màxima}) = 2\pi\nu A = 1,88\ \mathrm{m/s} \ \text{\punts{0,2}}\\
a_y(t) &= \frac{\mathrm{d}v_y}{\mathrm{d}t} = -4\pi^2\nu^2 A\cos(2\pi\nu t) \ \text{\punts{0,3}}\\
&\Rightarrow a_y(\text{màxima}) = 4\pi^2\nu^2 A = 3,55\cdot 10^{2}\ \mathrm{m/s^2} \ \text{\punts{0,2}}
\end{align*}
\end{solapartat}
